\section*{Broader impact}
\label{sec:impact}

The system itself is small (a 27-way classifier on cached features),
but a few of its design choices have implications worth being explicit
about. The default reference tree is built from a Western,
lineage-based art-history canon. Every metric in this paper that
mentions ``the hierarchy'' is anchored to that taxonomy, which is
implicitly endorsed by anyone using the numbers. Our tree-variant
experiments are partial mitigation: they document how much each
conclusion depends on the choice of tree. A deployment in a museum or
classroom should treat the tree as configuration, not as a constant.
WikiArt is heavily biased toward European painting; the corpus has
27 styles but only one (Ukiyo-e) sits outside the European-and-American
canon, and East Asian ink-painting traditions that span centuries are
collapsed into that single label. A retrieval system trained on this
data will under-rank non-Western works for ambiguous queries.
Hyperbolic geometry, which tightens local neighborhoods, could
compound the effect by making those already-tight neighborhoods more
confident. We therefore do not recommend deploying
hyperbolic style embeddings for attribution or authentication tasks
without expert human review: $50$--$65\%$ top-1 accuracy over $27$
well-known styles is far below the threshold any serious provenance,
insurance, or legal decision should require.
